\chapter{Needle Decompression of Pneumothorax}

\section*{Overview}
Needle decompression of pneumothorax is an emergency procedure in which a large-bore needle or catheter is inserted into the pleural space to release trapped air. It is the immediate, temporizing treatment for tension pneumothorax.

\section*{Alternative Names}
Needle thoracostomy, needle aspiration of pneumothorax, emergency pleural decompression

\section*{Principle}
In tension pneumothorax, a one-way valve effect allows air to enter the pleural space during inspiration but prevents its escape, progressively compressing the lung and mediastinum. Inserting a cannula into the pleural space converts the pressure to atmospheric, decompressing the lung and restoring venous return until a chest tube can be placed.

\section*{History}
Needle decompression has been a standard emergency intervention for tension pneumothorax since the mid-20th century and is taught in all advanced trauma and resuscitation courses. Its routine use outside tension physiology has been re-examined as simple aspiration for primary pneumothorax.

\section*{Why This Test or Procedure Is Ordered}
Performed emergently for suspected or confirmed tension pneumothorax with respiratory distress, hypotension, or cardiac arrest, and sometimes as simple aspiration for a large, symptomatic primary spontaneous pneumothorax.

\section*{Is This a Primary or Secondary Test or Procedure}
A primary emergency intervention for tension pneumothorax, and a secondary temporizing measure before definitive chest tube drainage in other settings.

\section*{How This Test or Procedure Is Performed}
The insertion site, typically the second intercostal space in the midclavicular line or the fourth or fifth intercostal space in the anterior axillary line, is identified and prepared. A large-bore over-the-needle catheter is inserted perpendicular to the chest wall just above the rib, a rush of air or blood confirms entry into the pleural space, and the catheter is advanced while the needle is withdrawn. The catheter is left open to air or connected to a one-way valve while definitive chest tube placement is arranged.

\section*{What Preparation Is Needed}
None is feasible in the emergency setting; sterile preparation of the chest wall and a readily available decompression kit are required. Chest radiography should not delay decompression in suspected tension pneumothorax.

\section*{Contraindications and Precautions}
There are essentially no absolute contraindications in a dying patient with tension pneumothorax; relative precautions include chest wall infection or coagulopathy. Care must be taken to avoid injury to the intercostal neurovascular bundle and to recognize that a single decompression may fail, especially in obese or muscular patients.

\section*{Risks}
Failure of decompression with recurrent tension, pneumothorax if no pneumothorax was present, bleeding, infection, lung laceration, and injury to mediastinal structures.

\section*{What Happens After the Test or Procedure}
The patient is reassessed for clinical improvement and undergoes chest radiography, followed by definitive tube thoracostomy when tension physiology is confirmed. The decompression catheter may be removed after a functioning chest tube is in place.

\section*{Understanding Your Results}
Clinical success is marked by immediate improvement in blood pressure, oxygenation, and breath sounds with a rush of air on entry. Radiographic confirmation shows lung re-expansion and correct catheter position before definitive drainage.

\section*{Normal Results}
Relief of tension physiology, re-expansion of the lung, and successful transition to chest tube drainage without recurrence.

\section*{Follow-up Tests and Procedures}
Chest radiography, arterial blood gas or pulse oximetry monitoring, and definitive chest tube placement are the standard follow-up. Repeat decompression may be required if the catheter occludes or the tension reaccumulates.

\section*{References}
\begin{enumerate}
  \item MedlinePlus. Pneumothorax. U.S. National Library of Medicine; 2025.
  \item Current Medical Diagnosis and Treatment. McGraw-Hill; 2025.
\end{enumerate}

\clearpage
